\chapter{Traceability Matrices}
\label{app:traceability-matrices}

This appendix provides a traceability matrix mapping the core architectural and algorithmic claims discussed throughout this monograph to their corresponding implementations across the five primary project workspaces (\texttt{wolverine-sih}, \texttt{wolverine-db}, \texttt{Vzeya}, etc.). This matrix ensures rigorous verifiability of the system's design principles by grounding theoretical claims in concrete source code artifacts.

\section{Core Algorithmic \& Architectural Claims}

\begin{table}[h!]
\centering
\caption{Traceability Matrix: Monograph Claims to Source Code Implementations}
\label{tab:traceability-matrix}
\renewcommand{\arraystretch}{1.6}
\begin{tabular}{|p{0.32\linewidth}|p{0.63\linewidth}|}
\hline
\textbf{Claim / Capability} & \textbf{Implementation Details \& Source References} \\
\hline
\textbf{30-Day Linear Temporal Decay} & Implementation of a strictly linear 30-day temporal decay function for entity confidence scores, explicitly rejecting exponential models. \\
& \textbf{Repository}: \texttt{wolverine-sih} \\
& \textbf{File}: \texttt{wolverine/src/resolver/entity\_resolver.ts} \\
\hline
\textbf{In-Memory Bounded BFS Graph Projection} & Graph traversal and entity resolution utilize purely in-memory bounded Breadth-First Search (BFS). PostgreSQL recursive CTEs are strictly avoided for graph projection purposes. \\
& \textbf{Repository}: \texttt{wolverine-sih} \\
& \textbf{File}: \texttt{wolverine/src/graph/projector.ts} \\
\hline
\textbf{CSS-Based CRT Shader Effects} & The terminal user interface (Vzeya) implements CRT visual effects (such as scanlines and vignettes) using CSS \texttt{repeating-linear-gradient} rather than WebGL GLSL shaders, despite directory naming. \\
& \textbf{Repository}: \texttt{Vzeya} \\
& \textbf{File}: \texttt{src/components/webgl/CRTPostProcessing.tsx} \\
\hline
\textbf{Simulated BFT Consensus} & Byzantine Fault Tolerance (BFT) consensus is implemented mathematically but operates over a simulated network using \texttt{DirectMemoryNetworkTransport}. \\
& \textbf{Repository}: \texttt{wolverine-db} \\
& \textbf{File}: \texttt{src/runtime/network\_transport.ts} \\
\hline
\textbf{0.92 and 0.70 Confidence Thresholds} & The exact threshold values used for entity resolution and linkage triage are strictly configured to 0.92 (high confidence auto-link) and 0.70 (moderate confidence manual review). \\
& \textbf{Repository}: \texttt{wolverine-sih} \\
& \textbf{File}: \texttt{wolverine/src/resolver/entity\_resolver.ts} \\
\hline
\end{tabular}
\end{table}

\section{Summary of Compliance}
The traceability matrix confirms that the implemented architecture consistently adheres to the constraints outlined in this monograph. Specifically, the strict avoidance of PostgreSQL recursive CTEs in favor of bounded in-memory BFS ensures predictable bounded-time resolution for analysts. Furthermore, the UI rendering optimizations utilizing CSS \texttt{repeating-linear-gradient} demonstrate adherence to the low-overhead terminal client specification. Simulated \texttt{DirectMemoryNetworkTransport} provides verifiable Byzantine tolerance without the operational overhead of a distributed physical network.
